\section{GraphDrone.fit() Implementation Program}

The implementation target is a portfolio-general GraphDrone object with four integrated layers:
\begin{enumerate}
\item a portfolio-level expert factory,
\item a per-row, per-view token builder,
\item a contextual set-router over variable expert sets,
\item a sparse defer-to-FULL integrator.
\end{enumerate}

This program is explicitly not a continuation of California-only or GEO-only tuning. The design is anchored in the registered signal-vs-noise portfolio, where specialist signal exists across multiple datasets but dense routing adds measurable competition noise.

For row $i$ and expert $v$, the target token is
\[
t_{i,v} = [p_{i,v}, q_{i,v}, s_{i,v}, d_v],
\]
where $p_{i,v}$ contains prediction-side information, $q_{i,v}$ contains prior-quality and competition-noise fields, $s_{i,v}$ contains support summaries, and $d_v$ is a learned or embedded descriptor for the expert family.

The contextual router produces sparse specialist scores $\alpha_{i,v}$ and a defer scalar $\delta_i$:
\[
\alpha_{i,v} = \mathrm{sparsemax}(W_\alpha h_i), \quad
\delta_i = \sigma(W_\delta h_i),
\]
with final prediction
\[
\hat y_i = (1 - \delta_i)\hat y_{i,\mathrm{FULL}} + \delta_i \sum_v \alpha_{i,v}\hat y_{i,v}.
\]

The implementation program is phase-gated. Phase I-A creates the package skeleton, portfolio loader, and public API. Phase I-B builds the expert factory. Phase I-C replaces scalar-only routing inputs with token and support-encoder tensors. Phase I-D implements the first lightweight contextual set-router. Phase I-E benchmarks the new object on the registered portfolio.

External review tightened the anti-drift constraints:
\begin{itemize}
\item no scalar-only routing surrogate,
\item no fixed-weight averaging used as a stand-in for contextual routing,
\item no benchmark scripts owning model internals,
\item no benchmark-path coupling inside the model package,
\item no silent collapse to FULL,
\item no dataset-specific assumptions such as mandatory GEO experts.
\end{itemize}

\begin{table}[h]
\centering
\begin{tabular}{p{0.12\linewidth} p{0.29\linewidth} p{0.22\linewidth} p{0.23\linewidth}}
\hline
Phase & Goal & Core Deliverables & Gate \\
\hline
I-A & Create the package skeleton, portfolio loader, and public API & \texttt{src/graphdrone\_fit/}, \texttt{portfolio\_loader.py}, \texttt{GraphDrone.fit()}, typed config objects & Model package exists without benchmark logic or benchmark-path coupling \\
I-B & Build the portfolio expert factory & expert ontology, descriptor validation, registry bridge & Variable expert sets instantiate without router hardcoding \\
I-C & Replace scalar routing inputs with tokens and support tensors & token schema, support encoder, descriptor embeddings & Every routing decision is traceable to row-level tokens \\
I-D & Implement the first contextual set-router & sparse router, defer gate, diagnostics & Beats fixed hedge on the registered portfolio without collapsing to dense averaging \\
I-E & Benchmark and log the new object & portfolio runner, leaderboard integration, milestone artifacts & Full comparison against registered runners and saved diagnostics \\
\hline
\end{tabular}
\caption{Phase-gated implementation program for the portfolio-general GraphDrone.fit() architecture.}
\end{table}

